\section{Readiness, interfaces, and process accounting}
\label{app:readiness}

\subsection{Prospective readiness measurement}

For a case class, enumerate required facts, decision criteria, evidence claims, actions, and
acceptance tests independently of the evaluated agent. Define
$d_i=(O_i,C_i,V_i,X_i,Y_i)\in[0,1]^5$ using Table~\ref{tab:readiness}. Undefined denominators are
reported as such. They do not become perfect scores.

\begin{table}[!htb]
\centering\small
\caption{Candidate readiness dimensions. Ratios are measured before evaluated outcomes.}
\label{tab:readiness}
\begin{tabularx}{\linewidth}{@{}l LL@{}}
\toprule
Score & Numerator / denominator & Distinction\\
\midrule
$O$ & Accessible required facts / all required facts. & Authorized accessibility, not truth.\\
$C$ & Explicit testable criteria / enumerated criteria. & Auditable requirements and ambiguity
handling, not a complete encoding of expert thought.\\
$V$ & Claims with independent checking routes / material claims. & Verification access, not a claim
that every check passed.\\
$X$ & Executable scoped actions / required actions. & Technical capability, not authority to approve.\\
$Y$ & Outcomes with independent tests or adjudication / required outcomes. & Evaluation coverage,
not observed success.\\
\bottomrule
\end{tabularx}
\end{table}

A candidate summary is $A_i=(O_iC_iV_iX_iY_i)^{1/5}$. It is a descriptive readiness index, not a
success probability. Its value for $(0.01,1,1,1,1)$ is about 0.398; a minimum index and individual
hard blockers provide useful alternative predictors. An unavailable mandatory action can block
deployment even if most scored actions are accessible. Institutional permission $\Gamma_i$ is
recorded separately. The formalization need not encode every step of tacit expert reasoning:
independent outcome adjudication can exist where procedural rules are incomplete.

For an artifact change $\Delta z_{ij}=\max(0,z^1_{ij}-z^0_{ij})$, a bounded positive-change
representation is
\begin{equation}
d'_{ja}=d_{ja}+(1-d_{ja})\beta_{ija}\Delta z_{ij},\qquad 0\le\beta_{ija}\le 1. \label{eq:readiness-update}
\end{equation}
Here $z_{ij}$ is normalized to $[0,1]$, so that $\Delta z_{ij}\in[0,1]$ and the updated
dimension stays in $[0,1]$. Only dimensions affected by the artifact receive nonzero sensitivity. The readiness index is
recomputed from the new vector. An unchanged artifact makes no change, an unaffected zero
dimension remains zero, and authority remains separate. Degradation is evaluated as its own
outcome rather than inferred from this positive-change representation. The labor response remains
Equation~\eqref{eq:response} and requires estimation.

\subsection{Visits, allocation, and baseline coherence}

For a stable expected-flow model, $\lambda=v+R^{\top}\lambda$, hence $\lambda=(I-R^{\top})^{-1}v$
when the spectral radius of $R$ is below one. $v$ counts exogenous visits and $R_{ij}$ expected
follow-on visits to $j$ from one visit to $i$. Branching and remediation can make visits exceed new
projects. Direct logs are preferred when an expected routing model is inappropriate.

The synthetic case partitions in Table~\ref{tab:partition} satisfy Equation~\eqref{eq:partition}.
The frozen partition separates baseline concentration from subsequent inflation; post-deployment
multipliers are $k=\kappa\eta$. The reproduction code checks this constraint, the four workload
totals, the fragility margins, and the route illustration.

\begin{table}[!htb]
\centering\small
\caption{Baseline partition, symmetric review effort, and fragility margin. $q\kappa=\phi$ is retained
baseline case labor. $\mu=h_R/h_S^H=m/\sigma$ measures review against its own group's baseline, with
$m=0.10$ in S0--S2 and $m=0.35$ in S3. $\eta^{*}$ is the exception-effort ratio at which workload
equals the baseline under the S2 review, rework, and upkeep (Equation~\ref{eq:etastar}).}
\label{tab:partition}
\begin{tabular}{@{}l r r r r r r r@{}}
\toprule
Pool & $q$ & $\kappa$ & $\sigma$ & $q\kappa$ & $\mu$ (S0--S2) & $\mu$ (S3) & $\eta^{*}$ (S2)\\
\midrule
\rowsinput{generated/table_partition_rows}
\bottomrule
\end{tabular}
\end{table}
